\documentclass[10pt,letterpaper]{article}
\usepackage{cornell-notes}

\title{Chapter 1: Next Subset of an n-Set}
\author{}
\date{}

\setDocTitle{Chapter 1: Next Subset of an n-Set}
\setDocSubtitle{Combinatorial Algorithms}
\setDocVersion{v1.0}
\setDocStatus{Study Companion}
\setDocOwner{Combinatorial Algorithms Cornell Notes}
\setDocAuthor{}
\setDocDate{\today}
\setDocSummary{Cornell-format study and review notes for Chapter 1: Next Subset of an n-Set.}

\setCornellCollection{Combinatorial Algorithms Cornell Notes}
\setCornellUnitType{Chapter}
\setCornellUnitNumber{1}
\setCornellUnitTitle{Next Subset of an n-Set}
\setCornellCompanionLabel{Study and review companion}

\hypersetup{pdftitle={Chapter 1: Next Subset of an n-Set},pdfsubject={Cornell notes},pdfauthor={}}

\begin{document}
\maketitle
\chaptermeta{1}{Next Subset of an $n$-Set (NEXSUB/LEXSUB)}{Sequential generation of all subsets}
\begin{tcolorbox}[colback=Teal!5,colframe=Teal,title=Essential question]
How should all $2^n$ subsets be sequenced when the surrounding computation benefits from direct binary order, one-element changes, or lexicographic ancestry?
\end{tcolorbox}
\begin{tcolorbox}[colback=white,colframe=Mist,title=Learning objectives]
\begin{itemize}
\item Encode a subset as a bit vector and generate its direct successor.
\item Explain why the reflected Gray code changes exactly one membership bit.
\item Apply the NEXSUB parity rule and maintain the subset cardinality.
\item Generate lexicographic successors with LEXSUB and explain its jump control.
\item Match each ordering to the external calculation it is intended to support.
\end{itemize}
\textbf{Prerequisites:} finite sets, binary representation, cardinality, lexicographic order, and basic graph terminology.
\end{tcolorbox}
\section*{Cornell notes}
\begin{longtable}{@{}Q!{\color{Mist}\vrule width 1pt}A@{}}
\cornellhead
\cue{What is the common problem?}{Generate every subset $S\subseteq\{1,\ldots,n\}$ once, returning one set per call while coordinating initialization and final termination. A bit vector $a_i\in\{0,1\}$ records membership and $k=|S|$ records cardinality. The chapter presents three orderings because the best successor rule depends on the computation $C(S)$ performed between outputs.}
\cue{How does direct binary order work?}{Associate $S$ with a binary integer whose bits are $a_1,\ldots,a_n$. The successor simulates addition of one: starting at $i=1$, turn consecutive $1$ bits into $0$ bits, decrementing $k$ for each deletion; then turn the first $0$ bit into $1$ and increment $k$. In set language, insert the smallest absent element and delete every smaller present element.}
\cue{What is the Gray-code goal?}{Consecutive subsets must differ by exactly one element. The bit vectors are vertices of the $n$-cube, so the desired sequence is a Hamilton walk through all vertices. Recursively, append $0$ to the $(n-1)$-dimensional list and append $1$ to its reversal. Reflection joins the halves by a single-bit transition and preserves single-bit changes inside both halves.}
\cue{What is NEXSUB's succession rule?}{For the current subset of cardinality $k$: if $k$ is even, change coordinate $j=1$; if $k$ is odd, scan from the first coordinate until the bit following the first $1$ is found, and change that coordinate. After toggling $a_j$, update $k\leftarrow k+2a_j-1$ using the new bit. The index $j$ identifies the sole changed coordinate.}
\cue{Why is the Gray rule correct?}{The reflected construction supplies an induction. The rule generates the first half exactly as it generates the lower-dimensional list. The appended final $1$ reverses cardinality parity in the second half, which reverses the lower-dimensional traversal. Thus the same local rule walks the reflected half in the required direction and visits each bit vector once.}
\cue{How is lexicographic order represented?}{Represent $S$ by its increasing member list $(a_1,\ldots,a_k)$. A proper extension of a prefix follows that prefix contiguously. Consequently, the most recent $(k-1)$-element predecessor of a $k$-set is its prefix, which lets an external computation retain results by depth in a stack.}
\cue{What does LEXSUB do?}{On initialization, extend the current prefix with its next legal element. For an ordinary successor, either append $a_k+1$ or, when the last element is $n$, back up until a position can advance. A logical jump control skips the entire contiguous block of supersets of the current set; this supports pruning when no extension can be useful.}
\cue{State and representation}{\begin{itemize}
\item \textbf{NEXSUB:} $n$, bit array, more-to-come flag, cardinality, and changed index.
\item \textbf{LEXSUB:} $n$, current length $k$, increasing member array, and jump flag.
\item Both routines mutate persistent state across calls; callers must not overwrite bookkeeping values between outputs.
\end{itemize}}
\cue{Complexity and cautions}{The source establishes bounded average work per generated subset for the direct and Gray-code rules; Gray transitions change one coordinate. LEXSUB may back up through a suffix. Guard the final state, keep bit and member-list representations distinct, and do not confuse lexicographic ordering of increasing lists with numeric bit-vector order.}
\cue{Retrieval practice}{\begin{enumerate}
\item What external computation favors Gray order? \item How is $k$ updated after a toggle? \item Why does reflection reverse the second half? \item What prefix property motivates LEXSUB? \item What does the jump flag discard? \item Which state must persist across calls?
\end{enumerate}}
\end{longtable}
\begin{tcolorbox}[breakable,colback=Ink!4,colframe=Ink,title=Cornell summary]
The chapter treats subset generation as an interface between an enumerator and a computation performed on every output. Direct binary order offers a simple increment-like rule. Reflected Gray order models subsets as vertices of a cube and supplies a Hamilton traversal in which exactly one membership bit changes; NEXSUB implements the traversal from the parity of the current cardinality and the position of the first set bit. Lexicographic order instead preserves prefix structure. LEXSUB exploits that structure to make the result for a subset's immediate prefix available and to skip all supersets of a rejected prefix. The central lesson is not that one order dominates, but that representation and succession order should be selected together with the incremental work required by $C(S)$. Correct use also depends on maintaining call-to-call state and recognizing the final output without duplicating or omitting the empty set.
\end{tcolorbox}
\end{document}
